
\documentclass[12pt,a4paper]{article}

\usepackage[margin=1in]{geometry}
\usepackage{amsmath,amssymb}
\usepackage{booktabs}
\usepackage{array}
\usepackage{enumitem}
\usepackage{xcolor}
\usepackage{hyperref}
\usepackage{fancyhdr}
\usepackage{setspace}

\hypersetup{colorlinks=true,linkcolor=blue,urlcolor=blue}
\pagestyle{fancy}
\fancyhf{}
\lhead{Statistical Learning}
\rhead{Exercise Answers}
\cfoot{\thepage}
\onehalfspacing

\title{\textbf{Statistical Learning Exercise Answers}\\
\large Chapter 1 and Linear Regression}
\author{Student Name: \underline{\hspace{6cm}}\\
Student ID: \underline{\hspace{6cm}}}
\date{\today}

\begin{document}
\maketitle

\section*{Reference}
These answers are based on \textit{An Introduction to Statistical Learning:
With Applications in Python (ISLP)} and the exercises provided for the
course. The ISLP book itself labels Linear Regression as Chapter 3, although
the course materials may refer to it as Chapter 2.

\section*{CHAPTER ONE}

\subsection*{Question 1: Flexible versus Inflexible Statistical Learning Methods}

\subsubsection*{(a) Very large sample size and small number of predictors}

\textbf{Answer: A flexible method is generally better.}

A very large sample provides enough observations to support a more complex
model and reduces the risk of overfitting. Since the number of predictors is
small, the model is also relatively easy to estimate.

\textbf{Memory point:} Large $n$ makes flexibility safer.

\subsubsection*{(b) Very large number of predictors and small number of observations}

\textbf{Answer: An inflexible method is generally better.}

When there are many predictors but few observations, a flexible method can
fit the training data too closely and overfit. A less flexible method is
generally more stable and has lower variance.

\textbf{Memory point:} Small $n$ + large $p$ means avoid excessive flexibility.

\subsubsection*{(c) Highly non-linear relationship between predictors and response}

\textbf{Answer: A flexible method is generally better.}

A flexible method can capture complicated and non-linear relationships
between the predictors and the response, whereas a highly inflexible model
may be too restrictive.

\textbf{Memory point:} Non-linear relationship means flexibility helps.

\subsubsection*{(d) Extremely high error variance}

\textbf{Answer: An inflexible method is generally better.}

High error variance means that the data contain substantial random noise.
A highly flexible method may attempt to fit this noise and overfit the
training data. A simpler model is generally more stable.

\textbf{Memory point:} High noise means keep the model simple.

\begin{center}
\begin{tabular}{l l l}
\toprule
Part & Better method & Main reason\\
\midrule
(a) & Flexible & Large sample reduces overfitting risk\\
(b) & Inflexible & Small sample and many predictors increase overfitting\\
(c) & Flexible & Captures non-linear relationships\\
(d) & Inflexible & High noise makes flexibility risky\\
\bottomrule
\end{tabular}
\end{center}

\newpage

\subsection*{Question 10: Boston Housing Data}

\subsubsection*{(a) Loading the Boston data set}

The Python version of ISLP uses the \texttt{ISLP} package. The data can be
loaded using:

\begin{verbatim}
from ISLP import load_data
Boston = load_data("Boston")
\end{verbatim}

\subsubsection*{(b) Number of rows and columns}

The Boston data set contains:

\[
\boxed{506\text{ rows and }13\text{ columns}}
\]

Each row represents a Boston suburb/census tract. The columns represent
characteristics of the suburb, including crime rate, zoning, industrial
activity, Charles River location, pollution, number of rooms, housing age,
distance to employment centres, highway accessibility, tax rate,
pupil-teacher ratio, lower-status population percentage, and median house
value.

\begin{center}
\begin{tabular}{ll}
\toprule
Variable & Meaning\\
\midrule
\texttt{crim} & Per-capita crime rate\\
\texttt{zn} & Proportion of residential land zoned for large lots\\
\texttt{indus} & Proportion of non-retail business land\\
\texttt{chas} & Charles River indicator\\
\texttt{nox} & Nitrogen oxide concentration\\
\texttt{rm} & Average rooms per dwelling\\
\texttt{age} & Proportion of older owner-occupied homes\\
\texttt{dis} & Distance to employment centres\\
\texttt{rad} & Accessibility to radial highways\\
\texttt{tax} & Property-tax rate\\
\texttt{ptratio} & Pupil-teacher ratio\\
\texttt{lstat} & Lower-status population percentage\\
\texttt{medv} & Median value of owner-occupied homes\\
\bottomrule
\end{tabular}
\end{center}

\subsubsection*{(c) Pairwise scatterplots}

The scatterplots show several relationships among the variables. In
particular, \texttt{rad} and \texttt{tax} have a strong positive relationship,
while \texttt{indus} and \texttt{nox} are positively related. \texttt{age}
and \texttt{nox} also tend to be positively related.

There are negative relationships between \texttt{dis} and several other
variables. In addition, \texttt{rm} and \texttt{lstat} show a negative
relationship. Overall, the scatterplots indicate that several predictors
are correlated with one another.

\subsubsection*{(d) Predictors associated with per-capita crime rate}

Yes. Crime rate is associated with several predictors.

The strongest positive associations include \texttt{rad}, \texttt{tax},
\texttt{indus}, \texttt{nox}, and \texttt{lstat}. Crime is negatively
associated with variables such as \texttt{dis}, \texttt{zn}, and
\texttt{medv}.

The crime-rate variable ranges from approximately
\[
0.00632\text{ to }88.9762,
\]
and its distribution is strongly right-skewed.

\subsubsection*{(e) High crime, tax rates, and pupil-teacher ratios}

Some suburbs have particularly high crime rates. The range of \texttt{crim}
is approximately
\[
0.00632\text{ to }88.9762.
\]

The tax rate ranges from
\[
187\text{ to }711.
\]

The pupil-teacher ratio ranges from
\[
12.6\text{ to }22.0.
\]

Crime has much more extreme variation than the pupil-teacher ratio. The
crime distribution is strongly right-skewed, with a small number of suburbs
having exceptionally high crime rates.

\subsubsection*{(f) Suburbs bordering the Charles River}

The variable \texttt{chas} equals 1 when a suburb bounds the Charles River
and 0 otherwise.

There are:
\[
\boxed{35\text{ suburbs}}
\]
that bound the Charles River.

\subsubsection*{(g) Median pupil-teacher ratio}

The median pupil-teacher ratio is:
\[
\boxed{19.05}
\]

Thus, the median is approximately 19.05 pupils per teacher.

\subsubsection*{(h) Lowest median value of owner-occupied homes}

Two suburbs have the lowest median house value: observations 399 and 406.
For both,
\[
\texttt{medv}=5,
\]
which corresponds to \$5,000 because \texttt{medv} is measured in thousands
of dollars.

These two observations have several unusual characteristics, including
very high crime, older housing, high highway accessibility, relatively high
taxes, and relatively high pupil-teacher ratios. They also have relatively
low numbers of rooms and relatively high values of \texttt{lstat}.

\textbf{Conclusion:} The two lowest-value observations have several
unfavourable socioeconomic and environmental characteristics.

\subsubsection*{(i) Suburbs with more than seven and eight rooms}

There are:
\[
\boxed{64}
\]
suburbs with more than seven rooms per dwelling.

There are:
\[
\boxed{13}
\]
suburbs with more than eight rooms per dwelling.

The 13 suburbs with more than eight rooms are relatively uncommon and tend
to have higher median house values and lower percentages of lower-status
population.

\newpage

\section*{LINEAR REGRESSION}

\subsection*{Question 1: Hypothesis Tests and p-values}

The question concerns the multiple linear regression model relating
\textit{Sales} to \textit{TV}, \textit{radio}, and \textit{newspaper}
advertising.

For each predictor, the null hypothesis is that the corresponding
advertising medium has no association with Sales after accounting for the
other advertising media.

\subsubsection*{TV}

\[
H_0:\text{TV advertising spending has no association with Sales,}
\]
after accounting for radio and newspaper advertising.

The p-value is less than 0.0001. Therefore, we reject $H_0$.

\textbf{Conclusion:} There is very strong evidence that TV advertising
spending is associated with Sales after accounting for radio and newspaper
advertising.

\subsubsection*{Radio}

\[
H_0:\text{radio advertising spending has no association with Sales,}
\]
after accounting for TV and newspaper advertising.

The p-value is less than 0.0001. Therefore, we reject $H_0$.

\textbf{Conclusion:} There is very strong evidence that radio advertising
spending is associated with Sales after accounting for TV and newspaper
advertising.

\subsubsection*{Newspaper}

\[
H_0:\text{newspaper advertising spending has no association with Sales,}
\]
after accounting for TV and radio advertising.

The p-value is 0.8599. Since this is greater than 0.05, we fail to reject
$H_0$.

\textbf{Conclusion:} There is insufficient evidence that newspaper
advertising spending is associated with Sales after accounting for TV and
radio advertising.

\begin{center}
\begin{tabular}{lcc}
\toprule
Predictor & p-value & Conclusion\\
\midrule
TV & $<0.0001$ & Significant association\\
Radio & $<0.0001$ & Significant association\\
Newspaper & $0.8599$ & Not significant\\
\bottomrule
\end{tabular}
\end{center}

\newpage

\subsection*{Question 6: Least Squares Line}

The fitted simple linear regression line is
\[
\hat{y}=\hat{\beta}_0+\hat{\beta}_1x.
\]

The least-squares estimate of the intercept is
\[
\hat{\beta}_0=\bar{y}-\hat{\beta}_1\bar{x}.
\]

Set $x=\bar{x}$:
\[
\hat{y}
=
\hat{\beta}_0+\hat{\beta}_1\bar{x}.
\]

Substituting the intercept estimate gives
\[
\hat{y}
=
(\bar{y}-\hat{\beta}_1\bar{x})
+\hat{\beta}_1\bar{x}
=
\bar{y}.
\]

Therefore, the least-squares regression line always passes through
\[
\boxed{(\bar{x},\bar{y})}.
\]

\textbf{Memory point:} At average $X$, predicted $Y$ equals average $Y$.

\newpage

\subsection*{Question 11: Regression Without an Intercept}

The exercise generates data according to
\[
x\sim N(0,1),\qquad y=2x+\epsilon,
\]
using the specified random seed.

\subsubsection*{(a) Regress $y$ onto $x$ without an intercept}

The model is
\[
Y=\beta X+\epsilon.
\]

Using the specified data, the results are approximately:

\begin{center}
\begin{tabular}{lc}
\toprule
Quantity & Result\\
\midrule
$\hat{\beta}$ & 1.9762\\
$SE(\hat{\beta})$ & 0.11695\\
$t$-statistic & 16.8984\\
$p$-value & $6.23\times10^{-31}$\\
\bottomrule
\end{tabular}
\end{center}

We test
\[
H_0:\beta=0.
\]

Because the p-value is extremely small, we reject $H_0$.

\textbf{Conclusion:} There is extremely strong evidence of an association
between $x$ and $y$. The estimated slope is approximately 1.98, close to
the true slope of 2.

\subsubsection*{(b) Regress $x$ onto $y$ without an intercept}

The reversed model is
\[
X=\beta Y+\epsilon.
\]

The results are approximately:

\begin{center}
\begin{tabular}{lc}
\toprule
Quantity & Result\\
\midrule
$\hat{\beta}$ & 0.37574\\
$SE(\hat{\beta})$ & 0.02224\\
$t$-statistic & 16.8984\\
$p$-value & $6.23\times10^{-31}$\\
\bottomrule
\end{tabular}
\end{center}

The coefficient changes, but the t-statistic and p-value remain the same.

\subsubsection*{(c) Relationship between (a) and (b)}

\begin{center}
\begin{tabular}{lcc}
\toprule
 & $Y$ on $X$ & $X$ on $Y$\\
\midrule
Coefficient & 1.9762 & 0.37574\\
Standard error & 0.11695 & 0.02224\\
$t$-statistic & 16.8984 & 16.8984\\
$p$-value & $6.23\times10^{-31}$ & $6.23\times10^{-31}$\\
\bottomrule
\end{tabular}
\end{center}

Thus,
\[
\boxed{t_{Y\mid X}=t_{X\mid Y}}
\]
and
\[
\boxed{p_{Y\mid X}=p_{X\mid Y}}.
\]

The direction of regression changes the coefficient because the variables
have different scales, but it does not change the strength of evidence
against the null hypothesis.

\subsubsection*{(d) Derivation of the t-statistic}

For regression of $Y$ onto $X$ without an intercept,
\[
\hat{\beta}
=
\frac{\sum_{i=1}^{n}x_i y_i}
{\sum_{i=1}^{n}x_i^2}.
\]

The standard error is
\[
SE(\hat{\beta})
=
\sqrt{
\frac{\sum_{i=1}^{n}(y_i-x_i\hat{\beta})^2}
{(n-1)\sum_{i=1}^{n}x_i^2}
}.
\]

Let
\[
A=\sum x_i y_i,\qquad B=\sum x_i^2.
\]

Then
\[
\hat{\beta}=\frac{A}{B}.
\]

Expanding the residual sum of squares gives
\[
\sum(y_i-x_i\hat{\beta})^2
=
\sum y_i^2-\frac{A^2}{B}.
\]

Therefore,
\[
SE(\hat{\beta})
=
\frac{\sqrt{B\sum y_i^2-A^2}}
{\sqrt{n-1}\,B}.
\]

Hence,
\[
\boxed{
t=
\frac{
\sqrt{n-1}\sum_{i=1}^{n}x_i y_i
}{
\sqrt{
\left(\sum_{i=1}^{n}x_i^2\right)
\left(\sum_{i=1}^{n}y_i^2\right)
-
\left(\sum_{i=1}^{n}x_i y_i\right)^2
}
}
}.
\]

\subsubsection*{(e) Equality of the t-statistics}

The formula above is symmetric in $x$ and $y$. Since
\[
\sum x_i y_i=\sum y_i x_i,
\]
switching $x$ and $y$ leaves both the numerator and denominator unchanged.

Therefore,
\[
\boxed{t_{Y\text{ on }X}=t_{X\text{ on }Y}}.
\]

\subsubsection*{(f) Regression with an intercept}

When an intercept is included, the t-statistic for
\[
H_0:\beta_1=0
\]
is also the same whether $Y$ is regressed on $X$ or $X$ is regressed on
$Y$.

In simple linear regression with an intercept,
\[
t=r\sqrt{\frac{n-2}{1-r^2}},
\]
where $r$ is the sample correlation.

Because
\[
\operatorname{Cor}(X,Y)=\operatorname{Cor}(Y,X),
\]
the t-statistic is identical in both regression directions.

\[
\boxed{t_{Y\text{ on }X}=t_{X\text{ on }Y}}
\]

\section*{Final Revision Points}

\begin{itemize}[leftmargin=*]
    \item Small p-value $\Rightarrow$ reject $H_0$.
    \item Large p-value $\Rightarrow$ fail to reject $H_0$.
    \item In simple linear regression, the least-squares line passes through
    $(\bar{x},\bar{y})$.
    \item Flexible methods are useful when there is enough data or a
    complicated non-linear relationship.
    \item Inflexible methods are safer when the sample is small or the data
    contain substantial noise.
    \item Reversing simple regression changes the coefficient, but the
    t-statistic for a zero-slope test remains the same.
\end{itemize}

\end{document}
